\documentclass[11pt,a4paper]{article}

% ---------------------------------------------------------
% PREAMBLE (stable, warning-free, Overleaf-safe)
% ---------------------------------------------------------
\usepackage[utf8]{inputenc}
\usepackage[T1]{fontenc}
\usepackage{lmodern}
\usepackage{amsmath, amssymb, amsfonts}
\usepackage{bm}
\usepackage{geometry}
\usepackage{graphicx}
\usepackage{hyperref}
\usepackage{cite}
\usepackage{authblk}
\usepackage{physics}
\usepackage{mathtools}

\geometry{margin=2.4cm}

\title{\textbf{Companion XLVIII: Minkowski Functionals in the Relaxation-Driven Cyclic Cosmology}}
\author{Michael Lehmann \\ \texttt{mi.lehmann@gmx.de}}
\date{April 2026}

\begin{document}
\maketitle

\begin{abstract}
Minkowski functionals provide a complete morphological and topological description 
of random fields and have become a powerful tool for analyzing the large-scale 
structure of the Universe. In the standard cosmological model, the topology of the 
density field is determined by Gaussian initial conditions and non-linear 
gravitational evolution. In the Relaxation-Driven Cyclic Cosmology (RDCC), the 
scale-dependent evolution of the gravitational potential modifies the morphology, 
connectivity, and topology of the cosmic web. This Companion develops a continuous 
description of Minkowski functionals in RDCC, deriving how the relaxation scale 
influences the volume fraction, surface area, mean curvature, and Euler 
characteristic of the density field. The resulting framework provides a 
distinctive topological signature of the RDCC gravitational field.
\end{abstract}
% ---------------------------------------------------------
\section{Introduction}

The large-scale structure of the Universe is not only characterized by its power 
spectrum or bispectrum but also by its topology and morphology. Minkowski 
functionals provide a complete set of geometric descriptors that quantify the 
shape, connectivity, and curvature of isodensity surfaces. They capture 
information beyond correlation functions and are sensitive to non-Gaussianity, 
mode coupling, and the global structure of the cosmic web.

In RDCC, the gravitational potential evolves differently. The relaxation mechanism 
suppresses the decay of small-scale modes and enhances the coherence of large-scale 
modes. This modifies the topology of the density field in a scale-dependent manner, 
producing distinctive signatures in the Minkowski functionals. These signatures 
provide a direct probe of the relaxation scale $k_{\mathrm{rel}}$.
% ---------------------------------------------------------
\section{Minkowski Functionals of the Density Field}

For a three-dimensional field, the Minkowski functionals are defined as:

\begin{align}
    V_{0}(\nu) &= \text{Volume fraction}, \\
    V_{1}(\nu) &= \text{Surface area}, \\
    V_{2}(\nu) &= \text{Integrated mean curvature}, \\
    V_{3}(\nu) &= \text{Euler characteristic}.
\end{align}

Here $\nu$ is the density threshold in units of the standard deviation.

In the standard cosmological model, these functionals follow analytic expressions 
for Gaussian random fields, with deviations arising from non-linear evolution.

In RDCC, the scale-dependent evolution of the gravitational potential modifies the 
density field in a way that alters all four functionals.
% ---------------------------------------------------------
\section{RDCC Modifications to the Density Field}

The density contrast evolves according to
\begin{equation}
    \delta_{\mathrm{RDCC}}(k,a)
    = \delta(k,a)
      \left[
        1 - \chi_{\delta}
        \left(\frac{k}{k_{\mathrm{rel}}}\right)^{\alpha}
      \right].
\end{equation}

This modification suppresses small-scale fluctuations and enhances large-scale 
coherence. The resulting density field is smoother, with fewer small-scale peaks 
and a more coherent filamentary structure.

These changes propagate directly into the Minkowski functionals.
% ---------------------------------------------------------
\section{RDCC Minkowski Functionals}

The Minkowski functionals of a weakly non-Gaussian field can be written as
\begin{equation}
    V_{i}(\nu)
    = V_{i}^{(G)}(\nu)
      + \Delta V_{i}(\nu),
\end{equation}
where $V_{i}^{(G)}$ is the Gaussian prediction and $\Delta V_{i}$ encodes 
non-Gaussian corrections.

In RDCC, the corrections become
\begin{equation}
    \Delta V_{i,\mathrm{RDCC}}(\nu)
    = \Delta V_{i}(\nu)
      \left[
        1 - \chi_{V_{i}}
        \left(\frac{k}{k_{\mathrm{rel}}}\right)^{\alpha}
      \right].
\end{equation}

This modification produces:

- a smoother volume fraction curve,  
- reduced surface area at high thresholds,  
- enhanced mean curvature at intermediate thresholds,  
- and a distinctive shift in the Euler characteristic.

The Euler characteristic, in particular, becomes a sensitive probe of the 
relaxation scale.
% ---------------------------------------------------------
\section{Topology of the Cosmic Web in RDCC}

The Euler characteristic counts the number of connected components minus the number 
of tunnels plus the number of cavities. In RDCC, the suppression of small-scale 
modes reduces the number of isolated peaks and tunnels, while the enhanced 
coherence of large-scale modes increases the size and connectivity of filaments.

The topology of the cosmic web becomes smoother and more coherent, with fewer 
small-scale structures and more prominent large-scale features.

This topological signature provides a direct observational probe of the RDCC 
gravitational field.
% ---------------------------------------------------------
\section{Conclusion}

Minkowski functionals in the Relaxation-Driven Cyclic Cosmology reflect the 
scale-dependent evolution of the gravitational potential. The relaxation mechanism 
modifies the morphology, connectivity, and topology of the cosmic density field, 
producing distinctive signatures in the volume fraction, surface area, mean 
curvature, and Euler characteristic. These effects provide a powerful topological 
probe of the relaxation scale and the temporal structure of the RDCC gravitational 
field. The present Companion establishes the theoretical foundation for future 
studies of cosmic topology, non-Gaussianity, and the geometry of large-scale 
structure in RDCC.

% ========================
% References
% ========================

\begin{thebibliography}{10}

\bibitem{Flagship} Lehmann, M. (2026). Relaxation-Driven Cyclic Cosmology (RDCC): A Minimal CPT-Symmetric Framework with a Single Parameter. Flagship Paper. Zenodo: \url{https://zenodo.org/records/19744958} (v43.0 or later).

\bibitem{Zenodo} The complete RDCC companion series (I--LVII) is available at the same Zenodo record above.

% Thematisch relevante Companions
\bibitem{CompXVI} Companion XVI: Boltzmann Code and Modified Hierarchy.
\bibitem{CompXXXI} Companion XXXI: RDCC Halo Model and Non-Linear Structure.
\bibitem{CompXXXII} Companion XXXII--XXXIX: Galaxy--Halo Connection, Cosmic Web, Lensing, RSD, ISW, tSZ, Rotation Curves, CGM.

% Wichtige externe Referenzen
\bibitem{Planck2020} Planck Collaboration (2020), Astron. Astrophys. 641, A6.
\bibitem{DESI2024} DESI Collaboration (2024/2025), arXiv: [aktuelle $f\sigma_8$/BAO-Paper].
\bibitem{JWST2024} Maiolino et al. (2024), Nature (JWST high-$z$ galaxies/quasars).
\bibitem{Kaiser1987} Kaiser, N. (1987), Mon. Not. R. Astron. Soc. 227, 1 (RSD).
\bibitem{Bartelmann2001} Bartelmann, M. \& Schneider, P. (2001), Phys. Rep. 340, 291 (Weak Lensing Review).
\bibitem{HuOkamoto2002} Hu, W. \& Okamoto, T. (2002), Astrophys. J. 574, 566 (CMB Lensing).
\bibitem{LewisChallinor2006} Lewis, A. \& Challinor, A. (2006), Phys. Rep. 429, 1 (CMB Lensing Review).
\bibitem{NFW1997} Navarro, J. F., Frenk, C. S. \& White, S. D. M. (1997), Astrophys. J. 490, 493 (Halo Profiles).

\end{thebibliography}

\end{document}
